% Executive summary, one page. Every quantity is a macro from
% canonical_numbers.tex, so the text cannot drift from the analysis.
\section*{Summary}

Austin calls itself the Live Music Capital of the World, and both the city and
the state spend public money on that claim. This study asks a narrower question
than the slogan: what can be measured about how the people who work in music are
actually doing, and how do the programs aimed at them compare with what other
states and cities run?

The city and the state have both built real institutions for music. Those
institutions are small, unevenly administered, and pointed more at venues and
tourism than at what musicians earn. The people they exist to serve are also
largely invisible in the statistics governments use to make decisions.

\vspace{5pt}
\noindent\stattile{0.235\textwidth}{\NPumsAustinMusMedpernpPos}{Median earnings from all work, Austin musicians, 2025 dollars}
\stattile{0.235\textwidth}{\NPumsAustinMusSelfemp}{Of Austin musicians work for themselves, against \NPumsAustinAllworkSelfemp\ of all Austin workers}
\stattile{0.235\textwidth}{\NCityLmfShareOfHotFyTwoFive}{Of Austin hotel-tax revenue reaches the Live Music Fund}
\stattile{0.235\textwidth}{\NVenueCorePctTwoZeroOneNineTwoZeroTwoFive}{Real change in receipts at core live-music venues, 2019 to 2025}
\par\vspace{5pt}
\srcnote{Sources for the four tiles: Austin musician earnings and self-employment, U.S.\ Census Bureau American Community Survey 2020--2024 five-year Public Use Microdata Sample for Texas, Musicians and Singers only (occupation code 2752), \NPumsAustinMusNPos\ Austin records pooled over five years, downloaded from \url{https://www2.census.gov/programs-surveys/acs/data/pums/2024/5-Year/}. Live Music Fund share of hotel-tax revenue, City of Austin open budget data on data.austintexas.gov, fiscal 2025 actuals. Real change in core live-music venue receipts, Texas Comptroller Mixed Beverage Gross Receipts (dataset \texttt{naix-2893} at \url{https://data.texas.gov/}) matched to \NVenuePanelIsCurated\ named Austin live-music rooms, deflated to 2025 dollars with the national consumer price index. Each source is developed in the section that uses it and listed in the technical appendix.}
\par\vspace{6pt}

\subsection*{Austin musicians earn about half what other Austin workers earn, and hours worked do not account for it}

Austin musicians report median earnings from work of \NPumsAustinMusMedpernpPos\
(\NPumsAustinMusMedpernpPosMoe) against \NPumsAustinAllworkMedpernpPos\ for all
Austin workers, both measured in the American Community Survey among employed
adults aged 18 to 64 with positive earnings (developed in Section~\ref{sec:earnings}). Their median usual week is
\NPumsAustinMusMedwkhp\ hours and their median year \NPumsAustinMusMedwkwn\
weeks, matching the Austin workforce on both. Across employed Texans, a model
that holds age, sex, education, hours, weeks and metro area fixed still places
music occupations \NPumsRegMusicPctgap\ below otherwise similar workers, so
schedule length is not where the gap comes from. That figure is an association
among the characteristics the survey records, not the cost of choosing music.
Comparing Austin with other large Texas metros requires a broader occupation
definition, because Musicians and Singers alone yield too few records for a
Houston or Dallas-Fort Worth cell. On the broader music-occupation definition
(Musicians and Singers plus Music Directors and Composers), Austin's median is
\NPumsMetroAustinMusicMedpernp, against \NPumsMetroDfwMusicMedpernp\ in
Dallas-Fort Worth and \NPumsMetroHoustonMusicMedpernp\ in Houston. Austin has
the lowest point estimate, but the Austin sample is small and the margins
overlap in both pairwise comparisons, so the ordering is not statistically firm;
the author makes no claim beyond that no comparison here supports the assumption
that Austin pays musicians better than the other large Texas metros.

\subsection*{Most of this workforce is self-employed, which is why the official numbers miss it}

The American Community Survey shows that \NPumsAustinMusSelfemp\ of Austin
musicians work for themselves, against \NPumsAustinAllworkSelfemp\ of all
Austin workers.\footnote{Class-of-worker categories 6 and 7 (self-employed,
unincorporated and incorporated), Musicians and Singers only (occupation code
2752), employed 18--64 with positive earnings, ACS 2020--2024 five-year PUMS.
See Section~\ref{sec:payroll} for the composition of that self-employed group
and its trend since 2005.} The Bureau of Labor Statistics Occupational
Employment and Wage Statistics program builds its wage series from employer
reports rather than household responses, so it can see only payroll
jobs.\footnote{OEWS is a semi-annual establishment survey administered by BLS
and the state workforce agencies; establishments report the wage of every
worker on their payroll but report nothing for the self-employed, who fall
outside the sampling frame. Program documentation:
\url{https://www.bls.gov/oes/oes_ques.htm}. BLS also does not publish an
\emph{annual} musician wage for any city or year, because the work is too
irregular to annualise from an hourly rate.} That survey records
\NEmpPayrollMusiciansAustinLast\ musician jobs in the Austin metro in 2025, a
count much smaller than the household survey implies because self-employed
musicians never enter the payroll survey at all (see Section~\ref{sec:payroll}). The City of Austin bought a
third series, from the Creative Vitality Suite database, which combines federal
wage data with self-employment records and so does count self-employed
earnings.\footnote{City of Austin open data,
\textit{Median Earnings for Creative Occupations}, dataset
\texttt{2qxc-8cme} at \url{https://data.austintexas.gov/}, retrieved 1 August
2026. Creative Vitality Suite is maintained by Employment and Cultural Vitality
Data (EMSI). The city's series covers 2016 and 2017 only.} That series put 2017
musician pay at \NColAnchorCityWageVsOewsTwoZeroOneSeven\ of the federal figure
for the same occupation and year, direct evidence of what payroll data miss (see Section~\ref{sec:measurement}).

\subsection*{Live-music rooms are being lost; the ones that stay open trade like the market around them}

The Texas Comptroller receives a monthly filing from every business selling
liquor, reporting the location's gross mixed-beverage receipts. This study
tracks the subset of those filings that come from named Austin live-music
rooms.\footnote{The venue panel comprises \NVenuePanelIsCurated\ Austin
live-music rooms tracked monthly from January 2007 through June 2026, matched
to Comptroller filings by taxpayer name and location address, and deflated to
2025 dollars with the national consumer price index. It is a \emph{curated} set
rather than a census, so it over-represents established rooms in the central
city and excludes venues that hold no liquor permit. Receipts measure alcohol
sales at each room, not attendance, ticket revenue or what performers were
paid. Filings by festival companies operating under venue-linked names are
excluded from the panel; leaving them in would make one room appear to grow
more than thirteenfold. The panel and its construction are described in
Section~\ref{sec:venues} and the technical appendix.} Between 2019 and 2025,
receipts changed \NVenueCorePctTwoZeroOneNineTwoZeroTwoFive\ at the panel's
core live-music tier (rooms whose main business is live performance) and
\NVenueCitywidePctTwoZeroOneNineTwoZeroTwoFive\ across the full citywide
mixed-beverage universe. The two totals do not compare room for room, because
the city added
\NVenueCitywidePermitsPctTwoZeroOneNineTwoZeroTwoFive\ more permit locations
over the period; the tracked rooms held
\NVenuePanelShareTwoZeroOneNine\ of citywide receipts in 2019 and
\NVenuePanelShareTwoZeroTwoFive\ in 2025. Once each venue's own trend and
seasonality are absorbed, live-music rooms and other bars trade
indistinguishably on log monthly receipts after March 2020, an average
difference of \NVenueDidPct\ with a 95 percent interval running from about 14
percent below to about 17 percent above the rest of the market
(Section~\ref{sec:venues}). The aggregate decline is therefore explained by
which rooms exist, not by weaker per-room trade among survivors. What changed is which
rooms exist: \NVenueExitsTwoZeroTwoZeroTwoZeroTwoSixCore\ exits from the tracked
panel from 2020 onward were core live-music rooms, and the panel's net additions
since 2016 come to \NVenueNetAddsTwoZeroOneSixTwoZeroTwoSix\ rooms against
\NVenueNetAddsTwoZeroZeroSevenTwoZeroOneFive\ over the previous nine years, a
comparison Section~\ref{sec:venues} qualifies because the earlier figure counts
every room already trading when the receipts file opens in 2007.

\subsection*{Census respondents expect to stay in music; renters are far less sure they will stay in Austin}

The 2022 Greater Austin Music Census surveyed \NCensusNRespondents\ people
working in and around Austin's music economy: musicians and other music
creatives (\NCensusNSectorCreative\ respondents), music-industry workers such
as managers, technicians and studio staff (\NCensusNSectorIndustry), and venue
operators or independent promoters (\NCensusNSectorPresenter), plus a smaller
group outside those three sectors.\footnote{Sound Music Cities and City of
Austin Music \& Entertainment Division, \textit{Greater Austin Music Census},
2022. The instrument was distributed through music-community mailing lists and
city channels between 15 July and 12 September 2022. Respondent-level microdata
downloaded from City of Austin open data, dataset \texttt{an3p-3yqx}, retrieved
1 August 2026. It is a self-selected online sample rather than a probability
sample of Austin's music workforce, so its shares describe the people who
answered and no margin of error can be computed for any of them; a
probability-sample benchmark for the same population does not exist. This study
treats the census as the best available direct measurement of the workforce's
non-earnings answers, with that caveat carried into every figure it produces.}
Of \NCensusRetentionDenom\ respondents who answered both three-year questions,
\NCensusIntendContinueMusic\ expected to still be working in music in three
years and \NCensusIntendStayAustin\ expected to still live in greater Austin.
Housing tenure separates the Austin-retention answer more sharply than any other
characteristic tested:
\NCensusStayHousingTenureHomeowner\ of homeowners expected to stay against
\NCensusStayHousingTenureRenter\ of renters, with years of experience the next
strongest discriminant (Section~\ref{sec:census}). Most respondents who did not
say yes to Austin answered ``unsure'' rather than ``no.''

\subsection*{Austin and Texas have built real music programs; the record shows how small and uneven they are}

The left column is what the record shows has been built, the right what the same
record shows about size and reach.

\vspace{3pt}
{\footnotesize
\setlength{\tabcolsep}{4pt}
\renewcommand{\arraystretch}{1.0}
\noindent\begin{tabular}{@{}P{0.485\textwidth} P{0.485\textwidth}@{}}
\toprule
\textbf{What is working} & \textbf{What is not} \\
\midrule
Austin runs a dedicated public fund for musicians, financed by a fixed share of
hotel-tax revenue rather than by an annual appropriation.
& The fund is \NCityLmfShareOfHotFyTwoFive\ of city hotel-tax revenue, and
historic preservation draws \NCityHpToLmfRatioFyTwoFive\ as much from the same
tax, because it draws on hotel-tax money beyond the equal share of the 2019
increase the ordinance gave each fund (see Section~\ref{sec:city}). \\
\rowrule
Texas is the only state compared here that subsidises operating music venues,
and in the one rebate round with a published total it awarded its full annual
dedication.
& Through the same state office, Texas has dedicated moving-image money against
music money for the 2026--27 biennium at \NTxFilmMusicRatio\ (see Section~\ref{sec:state}). \\
\rowrule
Federal shuttered-venue grants brought \NSvogAustinTotal\ in nominal award
dollars to Austin-area recipients, a one-off federal program in which motion
picture theatre operators took slightly more than live venues did (see Section~\ref{sec:state}).
& Neither the 2017 extended-hours pilot nor the 2024 venue grants left a
detectable mark on venue receipts (see Section~\ref{sec:venues}). \\
\rowrule
Health coverage among Austin musicians is \NPumsAustinMusHicov\
(\NPumsAustinMusHicovMoe), indistinguishable within that margin from
\NPumsAustinAllworkHicov\ for all Austin workers.
& Neither the city nor the state provides that coverage, and private coverage
runs \NPumsAustinMusPrivcov\ for musicians against
\NPumsAustinAllworkPrivcov\ for the wider Austin workforce (see Section~\ref{sec:earnings}). \\
\rowrule
The city publishes an awardee list for each Live Music Fund cycle.
& What a musician could receive depended heavily on the year they applied, and
one cycle was skipped entirely (see Section~\ref{sec:city}). \\
\rowrule
Austin rents and home values have fallen further from their peaks than in any
comparison metro (see Section~\ref{sec:costs}).
& Austin's own music census has not asked a musician what they earn since the
2014 fielding, whose income questions covered calendar 2013 (see Section~\ref{sec:measurement}). \\
\bottomrule
\end{tabular}
}

\vspace{6pt}
{\footnotesize\itshape The state Comptroller records every venue's alcohol sales
and the city records what it spends. No public source records what a musician is
paid for a night's work, the thing the argument is actually about. This study
takes no position in that argument.}
